% The complete pulpit example deck: ordinary Beamer slides followed by
% the media-overlay demonstrations from media.tex.
%
% Build with `make -C examples combined.pdf`.

\documentclass[aspectratio=169]{beamer}

\usepackage[T1]{fontenc}
\usepackage{graphicx}
\usepackage{booktabs}

\hypersetup{
  pdftitle={pulpit example deck},
  pdfauthor={pulpit}
}

% A poster wrapped in its own link: the link rectangle and the visible poster
% coincide by construction, so the overlay lands exactly where the poster is.
\begingroup
\catcode`\&=12\relax
\gdef\defrunlink#1#2{\gdef#1{#2}}
\defrunlink\gifuri{run:media-assets/bouncing.gif?autostart&loop}
\defrunlink\videouri{run:media-assets/clip.mp4?autostart&mute}
\defrunlink\loopvideouri{run:media-assets/clip.mp4?loop&mute}
\defrunlink\htmluri{run:media-assets/bouncing-balls.html}
\endgroup

\newcommand{\overlay}[3][0.62\textwidth]{%
  \href{#2}{\includegraphics[width=#1]{#3}}%
}

\newcommand{\uri}[1]{{\small\texttt{#1}}}

\title{pulpit example deck}
\subtitle{Presentation basics and media overlays}
\author{pulpit}
\date{}

\begin{document}

\begin{frame}
  \titlepage
\end{frame}

\section{Presentation basics}

\begin{frame}{One simple idea}
  \begin{itemize}
    \item Present a standard PDF slide deck.
    \item Keep the audience focused on the current slide.
    \item Move forward when you are ready.
  \end{itemize}
\end{frame}

\begin{frame}{A clear structure}
  \begin{enumerate}
    \item Introduce the topic.
    \item Develop the main argument.
    \item Finish with a memorable conclusion.
  \end{enumerate}
\end{frame}

\begin{frame}{Two useful views}
  \begin{columns}[T]
    \begin{column}{0.45\textwidth}
      \begin{block}{Presenter}
        Controls the pace and previews what comes next.
      \end{block}
    \end{column}
    \begin{column}{0.45\textwidth}
      \begin{block}{Audience}
        Sees a clean, focused version of the current slide.
      \end{block}
    \end{column}
  \end{columns}
\end{frame}

\begin{frame}{Keep it readable}
  \begin{itemize}
    \item Use short headings.
    \item Prefer a few strong points per slide.
    \item Leave enough space for every idea to breathe.
  \end{itemize}

  \begin{alertblock}{Remember}
    Slides support the presentation; they are not the presentation.
  \end{alertblock}
\end{frame}

\begin{frame}{Ready to present}
  \centering
  \Large Open the PDF, check the display, and begin.
\end{frame}

\section{Media overlays}

\begin{frame}{How a deck declares media}
  A media overlay is a link around a poster:

  \begin{center}
    \uri{\textbackslash href\{run:clip.mp4?autostart\}\{\textbackslash includegraphics\{poster\}\}}
  \end{center}

  \texttt{run:} is the convention pdfpc and Impressive already read, so a deck
  written for them needs no changes. The link rectangle is the overlay region,
  and the poster is what every other PDF reader shows.
\end{frame}

\begin{frame}{Animated GIF}
  The poster is the GIF's own first frame; pulpit replaces it with the
  running animation.

  \begin{center}
    \overlay{\gifuri}{media-assets/bouncing-still.png}

    \uri{run:media-assets/bouncing.gif?autostart\&loop}
  \end{center}
\end{frame}

\begin{frame}{Video}
  The clip sits beside the document, which keeps the PDF small and the asset
  editable without rebuilding the deck.

  \begin{center}
    \overlay{\videouri}{media-assets/poster.png}

    \uri{run:media-assets/clip.mp4?autostart\&mute}
  \end{center}
\end{frame}

\begin{frame}{Interactive HTML}
  The same convention extends to a page: the HTML file's own directory is
  served to it, so its stylesheet and script resolve exactly as on disk. On
  the slide it becomes live and clickable.

  \begin{center}
    \overlay[0.58\textwidth]{\htmluri}{media-assets/balls-poster.png}

    \uri{run:media-assets/bouncing-balls.html}
  \end{center}
\end{frame}

\begin{frame}{Incremental reveal continuity}
  Each reveal step repeats the identical link, and pulpit collapses the
  repetitions into one overlay --- so the video keeps playing rather than
  restarting as the bullets appear.

  \begin{columns}[T]
    \begin{column}{0.48\textwidth}
      \overlay[\textwidth]{\loopvideouri}{media-assets/poster.png}
    \end{column}
    \begin{column}{0.46\textwidth}
      \begin{itemize}
        \item The video starts on the first step.
        \pause
        \item It does not restart here.
        \pause
        \item Nor here: one overlay, one session.
      \end{itemize}
    \end{column}
  \end{columns}
\end{frame}

\begin{frame}{What pulpit reads}
  \begin{center}
    \small
    \begin{tabular}{ll}
      \toprule
      Written in the deck & What it declares \\
      \midrule
      \texttt{run:clip.mp4} & video beside the document \\
      \texttt{run:spin.gif} & animated image beside the document \\
      \texttt{run:page.html} & interactive HTML \\
      \midrule
      \texttt{?autostart} & start on commit (also spelt \texttt{?autoplay}) \\
      \texttt{?loop} & repeat \\
      \texttt{?mute} & silence (video) \\
      \texttt{?start=12.5} & seek, in seconds (video) \\
      \bottomrule
    \end{tabular}

    \vspace{1em}
    A \texttt{run:} link to anything that is not a media file is never
    executed and never becomes an overlay.
  \end{center}
\end{frame}

\end{document}
